Put the large centres at $(\pm12,0)$. The small centres are $(0,\pm4\sqrt3)$ because their distance to either large centre is $8\sqrt3$. The centre quadrilateral has area $96\sqrt3$. Its angles are $60^\circ$ at the large centres and $120^\circ$ at the small centres. Subtracting the four sectors gives $96\sqrt3-27\pi-(\dfrac{2\pi}{3})(8\sqrt3-9)^2=96\sqrt3+(96\sqrt3-209)\pi$. Thus $a+b+c=\boxed{401}$.
\begin{asy}
import graph;import geometry;usepackage("amsmath");size(7cm);
real r=8*sqrt(3)-9;pair L=(-12,0),R=(12,0),U=(0,4*sqrt(3)),V=(0,-4*sqrt(3));draw(circle(L,9));draw(circle(R,9));draw(circle(U,r));draw(circle(V,r));draw(L--U--R--V--cycle,dashed+blue);label("$9$",L);label("$9$",R);label("$8\sqrt3-9$",U);label("$8\sqrt3-9$",V);
\end{asy}

